\documentclass[11pt]{amsart}
\usepackage{geometry}                % See geometry.pdf to learn the layout options. There are lots.
\geometry{letterpaper}                   % ... or a4paper or a5paper or ... 
%\geometry{landscape}                % Activate for for rotated page geometry
%\usepackage[parfill]{parskip}    % Activate to begin paragraphs with an empty line rather than an indent
\usepackage{graphicx}
\usepackage{amssymb}
\usepackage{epstopdf}
\DeclareGraphicsRule{.tif}{png}{.png}{`convert #1 `dirname #1`/`basename #1 .tif`.png}

\title{Ejemplo 5.3 }
\author{Enrique Zuñiga Zuñiga}
%\date{}                                           % Activate to display a given date or no date

\begin{document}
\maketitle
%\section{}
\noindent Sea DOMICILIO un registro formado por cuatro campos, uno de los cuales es numérico y los tres restantes son del tipo cadena de caracteres. Su representación queda como se muestra en la figura 5.2.  \newline\newline
DOMICILIO = REGISTRO  \newline
calle: cadena de caracteres   \newline
número: entero  \newline
ciudad: cadena de caracteres  \newline
país: cadena de caracteres \newline
{Fin de la definición del registro DOMICILIO}  \newline
$$
\begin{array}{|c|c|c|c|}
\hline
Calle           & Numero  & Ciudad  & Pais \\
\hline
  &   &  &  \\
\hline
\end{array}
$$
%\subsection{}



\end{document}  